% ---------------------------------------------------------------------------
\section{The writing rules of the Voynich manuscript}
\label{sec:rules}

This section describes how the text is written, from the word to the book. For each property we give the measurement,
its interval and the comparison texts, and we say whether the property was known. Section~\ref{sec:comparisons} gathers
the comparisons with languages, manuscripts and generators, and Section~\ref{sec:scorecard} turns the properties into
the scorecard of the generator.

\subsection{Word form and word boundaries}
\label{sec:boundaries}

The next glyph is very predictable: the conditional entropy $h_2$ of the running text (space included) is 2.23 bits,
against 2.54--4.65 in the 24 natural languages \citep[as found by][]{bennett1976,lindemannbowern2020}. Words follow a
rigid internal grammar \citep{stolfi2000,zattera2022}, and four properties follow from it:
\begin{itemize}
\item spaces are predictable from the two glyphs on either side (F1 0.86; higher than 23 of the 24 languages, the
  exception being Chinese in pinyin), and uncertain spaces fall where this rule is uncertain
  \citep[as observed by][]{feaster2020,rozanova2026};
\item cutting the glyph stream at other places allowed by the rule mostly yields words attested elsewhere (50\% of the
  wrongly cut pieces, against a median of 6.5\% in languages; compare the ``separable words'' of \citealt{zattera2022});
\item the vocabulary fills the most probable glyph sequences (43\% of them are Voynich words; median 5\% in languages);
\item frequent words are the most probable sequences (rank correlation 0.58 between form probability and frequency;
  median 0.09 in languages).
\end{itemize}
In version~1 we read these properties as evidence that spaces are weak boundaries inside a continuous chain of glyphs.
That reading does not hold. A null that keeps the word grammar and removes every link between words reproduces all four
(\expcode{e3c85}): with words drawn independently from the same vocabulary, or generated one by one from a glyph model
of order two trained on Voynich words, the F1 of the spaces is 0.88, the attested wrong cuts 43--45\%, the filling 42--43\%
and the frequency--form correlation 0.55--0.62. These properties show that word boundaries are marked redundantly by
the shape of the words, as \citet{stolfi2000}, \citet{feaster2020} and \citet{rozanova2026} had described, and as in
any script with positional forms. They do not show that spaces are weak.

\subsection{The junction between adjacent words}
\label{sec:junction}

\paragraph{A strong link} The last glyph of a word says a lot about the first glyph of the next: the junction is 0.188
bits beyond the in-line shuffle (95\% interval by bifolium 0.163--0.211). Per language the median is 0.097; two
languages of 24 exceed the Voynich, Sanskrit and Tagalog, which write phonological links between words
(\expcode{e3c84}). The volunteers' gibberish has 0.014, and published generators between 0 and 0.08, except
\emph{voynich-fingerprint} (0.16--0.17; \expcode{e3c93}), which chooses the start of each word from the last glyph of
the previous one. The link was
observed by \citet{currier1976}, \citet{smith2017lastfirst} and \citet{smithponzi2019}.

\paragraph{It mirrors the rules inside words} The preferences between the last glyph of a word and the first glyph of
the next resemble the preferences between consecutive glyphs inside a word: the rank correlation between the two tables
of pointwise mutual information is 0.47 (median of five samples of 10{,}000 words), against at most 0.37 in the 24
languages and 0.39 in the constructed ones; in the five samples drawn for the nulls of \expcode{e3c85} it is 0.55,
which shows how much the value varies between samples. This part does not follow from word form: on those samples
the nulls bring it down to 0.13 or less (0.22 when only the order of the words within each line is shuffled), and the
junction itself to zero. At equal $h_2$ the Voynich lies 4.6 standard deviations above the regression line of the
languages with the value 0.55, and 3.9 with 0.47. The correlation uses the cells seen at least five times in both
tables, so it leaves out the pairs that occur across spaces but hardly inside words (\eva{y.q}, \eva{n.o}), which
\citet{smithponzi2019} single out as marks of the boundary. \citet{smithponzi2019} concluded that word breaks are real boundaries, and
Section~\ref{sec:boundaries} agrees; what we add is that the dependency across the boundary follows the same pattern as
the dependency inside the word, more than in any language of the sample.

\paragraph{Two junction rules} (1) Before a gallows, a word begins with \eva{qo-} after a word ending in \eva{-y},
\eva{-o} or \eva{-d} (59--64\% of the time) and with \eva{o-} after \eva{-n}, \eva{-r}, \eva{-s} or \eva{-m} (\eva{qo-}
only 23--27\%); in part this was described by \citet{currier1976} and \citet{smith2017o}. (2) A word ends in \eva{-r}
before a word beginning with \eva{a-} (85\%) and in \eva{-l} before \eva{k-}, \eva{t-}, \eva{d-}, \eva{l-}, \eva{s-},
\eva{q-} (63--86\%). Both rules hold for each position in the line where there are enough cases
(\expcode{e3c87}; for \eva{-l}/\eva{-r} only the class of the fourth word onwards has enough, and the other classes go
the same way): with a null that keeps position, the junction stays at 0.196 against 0.188, and in each interior
position \eva{qo-} follows \eva{-y}/\eva{-o}/\eva{-d} 61--70\% of the time against 23--33\% after
\eva{-n}/\eva{-r}/\eva{-s}/\eva{-m}. The caution of \citet{feaster2022} does not apply to these rules. They go
the same way, with intervals above zero, in the three main hands (with a strength that varies up to twofold), in every
section with enough data, and in the three transliterations. They have the shape of
rules such as English \emph{a}/\emph{an} or French \emph{liaison}; they say nothing about the sounds behind them, and
they are compatible with a phonetically written language.

\paragraph{Looking one word ahead, and no further} When two words copied from the line above would jointly violate a
rule, the piece that changes is the end of the first word for \eva{-l}/\eva{-r} (59\% against 21\%) and the start of the
second for \eva{qo-}/\eva{o-} (60\% against 18\%). The pattern is compatible with writing that knows the next word while
finishing the current one; it rests on few conflicts (37 and 60) and does not survive the strictest multiplicity
correction. Beyond the next word there is no link: for a given middle word, the last glyph of a word says nothing about
the first glyph of the word after next (excess 0.003 bits, $z = 1.3$), while 22 of 24 languages have such a link.

\subsection{The line}
\label{sec:line}

\paragraph{The junction stops at the line break} Between the last word of a line and the first word of the next the link
is zero ($-0.003$ bits, $z = -0.7$), against 0.19 inside the line; the ratio $Q$ of the two is $-0.02$ on the whole text
and 0.02 on samples of the length of the comparison texts (\expcode{e3c89}). In the comparison languages, whose lines follow
the source editions and often end at a clause or a verse (Section~\ref{sec:data}), the median $Q$ is 0.49 and only
Hebrew falls below 0.10; because those lines are not physical, the languages are a weak reference for line
properties, and texts written in real lines are a better one. A herbal in verse, the
\emph{Macer floridus} printed one verse per line, has a closed line too ($Q = -0.004$). The line is therefore a unit of
the writing, as \citet{currier1976} and \citet{rozanova2026} found; a line that is a unit of content (a verse, an entry,
a recipe) would produce the same closure (Section~\ref{sec:discussion}). The preregistered criterion of an earlier
experiment (\expcode{e384}), set from short and noisy texts, read literally gave the opposite (``the junction carries
over the line break as in prose''); \expcode{e3c89}, preregistered with its own criterion, confirms the closed line.

\paragraph{And at the gap left by a drawing} Where a plant interrupts a line, the junction between the words on the two
sides drops from 0.167 to 0.026 bits, and the word after the gap begins like a word at the start of a line.

\paragraph{Line-edge forms} For the same rest of the word, at line end \eva{-m} gains 15 percentage points and \eva{-r}
loses 12 (\eva{dar} inside a line, \eva{dam} at its end); at line start \eva{s-} and \eva{y-} gain 9--10 points in place
of \eva{ch-} and \eva{k-} ($z$ from 13.6 to 26.7). In the last line of a paragraph, which does not reach the margin,
\eva{-m} is rarer: it is a form tied to the margin. These forms were described by \citet{currier1976},
\citet{smith2015linestart,smith2016m} and \citet{feaster2020,feaster2022}; we measure them for the same word and compare
them with scribes in Section~\ref{sec:scribes}.

\paragraph{Labels} In the labels of the drawings, which stand alone, \eva{qo-} before a gallows is almost absent (0--5\%
against 30--68\% in lines). We predicted this from the junction rule before looking at the labels; the fact itself had
been noticed by \citet{zandbergen_labels}.

\subsection{The lines above}
\label{sec:vertical}

\paragraph{Copying from the line immediately above} A word has an identical or nearly identical word in the two lines
above more often than the vocabulary of the paragraph predicts ($z = 17.4$), more than twice as much as in the generator
of \citet{timmschinner2020}, which was built on this observation of \citet{timm2014}. What we add concerns the shape of
the copying: it comes almost entirely from the line immediately above (in languages recurrence spreads over three or
four lines); two neighbouring copied words come from two neighbouring words above, usually in the same order; a copied
word takes the \eva{qo-}/\eva{o-} or \eva{-l}/\eva{-r} form that agrees with its new neighbour more than the form of its
source (effect $+0.49$ against $+0.16$); and copying stops at the paragraph and at the page.

\paragraph{Each line avoids starting like the line above} Shuffling the order of the lines within paragraphs shows that
a line avoids beginning like the line above: identical starts are 0.51 of the expected number (95\% interval by bifolium
0.40--0.61; 0.53 on samples of 10{,}000 words), with the strongest effects for \eva{qo-} ($z = -9.7$) and \eva{o-}
($-7.0$). If the line above starts with \eva{qo-}, the probability that the next starts with \eva{qo-} drops by 52
points. The effect concerns the left margin only: it is absent where lines resume to the right of a drawing, and the
line two above, the end of the line above and the right margin do not matter. In the languages the median ratio is 0.97
and one language of 23 falls below the Voynich (Portuguese; two when the text of each language closest to the
Voynich is used), but here too the languages, without physical lines, are a weak reference; none of the 79
manuscripts avoids the margin (Section~\ref{sec:scribes}); the volunteers' gibberish
has 1.32 and the generators 0.93--2.05. A forum study had reported vertical avoidance of the same glyph in the herbal
pages of one scribe \citep{tavie2024}; we measure it across the manuscript and against the comparison texts.

\subsection{Choices between variant forms}
\label{sec:choices}

Many Voynich words exist in pairs that differ by one glyph: \eva{qo-}/\eva{o-} at the start, \eva{k}/\eva{t} (two
gallows), \eva{sh}/\eva{ch}, \eva{-ey}/\eva{-dy} at the end. We treat these as \emph{variant forms}: choices whose
distribution we can measure for the same rest of the word. This is a working hypothesis, not a result. If the variants
were morphemes, their agreement between nearby words would be a kind of grammatical agreement, which languages have;
we therefore compare them both with grammatical agreement in languages and with allographs in medieval scribes.

\paragraph{A short state} Between nearby words of the same line, the same choice is repeated more than the preferences
of each word, page and position predict. With the bias of the measure corrected (Appendix~\ref{app:measures}), the
agreement of \eva{k}/\eva{t}, \eva{sh}/\eva{ch} and \eva{-ey}/\eva{-dy} together is $+0.132$ between adjacent words
(95\% interval by bifolium 0.111--0.153) and $+0.091$ at two and three words (0.069--0.113); it halves in about three
words. It is almost the same between very different words ($+0.09$) as between similar ones ($+0.13$), unlike the
copying generator of \citet{timmschinner2020}; it is not an effect of position in the line; it is separate for each
choice (agreement between different choices $-0.01$); and it is the same in hands 1, 2 and 3 ($+0.12$ to $+0.15$), in
languages A and B and in the three transliterations. Agreement of this size exists in languages as grammatical
agreement (Italian \emph{-o}/\emph{-a}, Latin \emph{-us}/\emph{-a}); the short state alone therefore does not set the
Voynich apart from a language with agreement. Its comparison with scribes is in Section~\ref{sec:scribes}. Local
clustering of similar forms was noted by \citet{timm2014} and \citet{schinner2007}, and a switch of the vowel after
\eva{ch}/\eva{sh} that holds for whole folios by \citet{parisel2026b}; we have not found the short, word-independent state
described before.

\paragraph{A slow state} Besides the short state there is a slow component that crosses the line break: words of
consecutive lines agree ($+0.05$), at two lines less ($+0.03$), at 7--12 lines not at all. It is not a trend from the
top to the bottom of the page, and it is normal in real scribes (Section~\ref{sec:scribes}).

\paragraph{Drift along the line} Inside the line, for the same word and leaving out the first and last word, the marked
variants \eva{qo}, \eva{k}, \eva{sh} and \eva{-ey} become rarer towards the right: from $-0.020$ to $-0.034$ per ten
glyphs, that is about 8--13 percentage points over a line of median length (39 glyphs). The drift is clearest for
\eva{sh}/\eva{ch} ($z = -6.3$ by bifolium); for the other three choices it has the same direction but each one alone does
not pass the strictest multiplicity correction. It is steeper in the first half of the line, present in every hand and
about three times stronger in language B. The rightward indices of \citet{feaster2021,feaster2022} (\eva{qo-} further left
than \eva{o-}, \eva{Sh} further left than \eva{ch}) describe the same tendency; we measure it for the same word and compare
it with scribes. Next to drawings the marked forms become rarer, and the \eva{q} drops in the word right after a drawing
($-33$ points). On the images, none of 16 sampled words in \eva{o-} after a drawing hides an untranscribed \eva{q}
(three are uncertain, where opaque paint touches the word), while the \eva{q} is plain in the four \eva{qo-} words
checked as a control: the drop is not, or not entirely, an artefact of transcription (\expcode{e3c94}).

\subsection{Repetition}
\label{sec:repetition}

One adjacent pair of words in a hundred is the same word repeated (\eva{chol chol}, \eva{qokeedy qokeedy}), and 3.9\% are
nearly identical (one glyph different), as reported by \citet{boxer2022} and \citet{bowernlindemann2021}. In version~1 we
read this as repetition beyond chance. The observed/expected ratio corrects that reading (\expcode{e3c88}):
\begin{itemize}
\item against the composition of each line (words shuffled within the line), the Voynich repeats the previous word as
  often as chance (ratio 1.04), whereas every comparison language and manuscript avoids it: languages 0.01--0.92 (the
highest is Sanskrit),
  the Nordic scribes 0.02--0.07, the German manuscripts at most 0.79, the Copiale copyist 0;
\item against the whole vocabulary ($\sum p^2$), the ratio is 2.85: repeated words gather in the same line (Sanskrit
  reaches 2.86).
\end{itemize}
The property of the Voynich is therefore that immediate repetition is \emph{not avoided} and that repeated words cluster
within lines; nearly identical neighbours are not distinctive (ratio 1.08, with four languages higher). Within the line,
repetition is strong between nearby words and fades after six or seven words, whereas in languages it grows with distance
(median ratio far/near 2.86); two languages of 20 fall below the Voynich, both with unstable negative ratios. When a
repeated word changes its start or end, the change follows the junction rule with its neighbour (\eva{okeey qokeey}).

\subsection{Paragraphs, pages and bifolia}
\label{sec:book}

\begin{itemize}
\item \textbf{The paragraph opens with a large sign}: 83\% of paragraph-initial words begin with a gallows, against 9\%
  elsewhere, and \eva{p} dominates (54\%) \citep{currier1976,dimperio1978,stafford2022,feaster2022}. The opening gallows
  is three times less tied to its word than usual: it is largely a mark of position.
\item \textbf{First and last lines have their own register}: longer words and more \eva{p}, \eva{f}, \eva{sh} and final
  \eva{-y} in the first line; fewer \eva{qo-} and gallows, more \eva{ch-}, \eva{s-}, \eva{y-} in the last.
\item \textbf{\eva{-ey} grows down the page}: for the same position in the line, the share of \eva{-ey} among words in
  \eva{-ey}/\eva{-dy} rises by 12 percentage points from top to bottom of the page (95\% interval 7--18); an apparent rise along
  the paragraph came from its opening lines. Lines get shorter down the page.
\item \textbf{Each page has a strong identity}, mainly along an axis from words in \eva{-edy}/\eva{-eey} to words in
  \eva{-aiin}/\eva{-ain}: pages depart from the profile of their section 2.8 times more than chance.
\item \textbf{The bifolium is a unit of identity}: the identity of a page is shared by the other side of the same leaf and
  by the conjugate leaf of the same bifolium, which in the bound book can lie far away, and not by the facing page
  ($z = 6.6$). The two halves of a bifolium share rare words and copy from each other. All 48 bifolia with at least two
  pages labelled for language are entirely in one language, and 94\% of the 52 with at least two pages labelled for
  hand are in one hand (these counts include pages without running text, hence 52 against the 50 bifolia of
Section~\ref{sec:data}). Hands and languages are recorded page by page in the files, and
  \citet{davis2025} had noted that in the herbal section the hands follow the bifolia; the shared rare words and the
  copying between conjugate leaves are what we add. Davis prefers, tentatively, quires written by one scribe and later
  bound out of order to scribes working on loose bifolia; our data do not separate the two.
\item \textbf{Languages A and B differ in vocabulary, not in glyphs}: the best substitution of glyphs brings A closer to B
  by only 1\% of the gap, while the same search recovers a random key every time. B is not A enciphered with another key.
\item Consecutive pages share more words than two pages of the same section, but mostly because they are the two sides of
  one leaf; this does not show that the present order of the leaves is the order of writing, and the binding is known not
  to be original \citep{davis2025}.
\end{itemize}
